\documentclass[9pt,a4paper]{extarticle}

\usepackage[utf8]{inputenc}
\usepackage[T1]{fontenc}
\usepackage[italian]{babel}
\usepackage{lmodern}
\usepackage{amsmath,amssymb}
\usepackage[a4paper,top=1.05cm,bottom=1.05cm,left=1.1cm,right=1.1cm]{geometry}
\usepackage{xcolor}
\usepackage{multicol}
\usepackage{enumitem}
\usepackage{microtype}

% ------------------------------------------------------------
% Impaginazione compatta
% ------------------------------------------------------------
\setlength{\parindent}{0pt}
\setlength{\parskip}{1.5pt}
\setlength{\columnsep}{0.55cm}
\setlength{\columnseprule}{0.2pt}

\setlength{\abovedisplayskip}{3pt}
\setlength{\belowdisplayskip}{3pt}
\setlength{\abovedisplayshortskip}{2pt}
\setlength{\belowdisplayshortskip}{2pt}

\pagestyle{empty}

% ------------------------------------------------------------
% Colori e comandi
% ------------------------------------------------------------
\definecolor{errore}{RGB}{150,30,30}
\definecolor{corretto}{RGB}{20,100,45}
\definecolor{titolo}{RGB}{40,60,105}

\newcommand{\R}{\mathbb{R}}

\newcommand{\err}[1]{%
	{\color{errore}\textbf{Errore:}} #1
}

\newcommand{\corr}[1]{%
	{\color{corretto}\textbf{Corretto:}} #1
}

\newcommand{\regola}[1]{%
	{\color{titolo}\textbf{Da ricordare:}} #1
}

\newcommand{\voce}[1]{%
	\vspace{3pt}
	{\color{titolo}\textbf{#1}}\par
	\vspace{1pt}
}

% ============================================================
\begin{document}

\begin{center}
	{\Large\bfseries Errori ricorrenti nel calcolo algebrico e nelle equazioni}\\[-1pt]
	{\small Classe 2BS -- Scheda di recupero e autocorrezione}
\end{center}

\vspace{-2mm}

{\small
Questa scheda raccoglie alcuni errori comparsi frequentemente.
L'obiettivo non è imparare divieti a memoria, ma capire
\textbf{che cosa stiamo facendo e perché}.
}

\vspace{1mm}

% ============================================================
% PAGINA 1
% ============================================================

\begin{multicols}{2}
\small

% ------------------------------------------------------------
\voce{1. Ridurre subito i termini simili}

Portarsi dietro termini già sommabili rende i calcoli più lunghi
e aumenta la probabilità di errore.

\err{
\[
3x-2+5x+7-4x
\]
e continuare così nei passaggi successivi.
}

\corr{
\[
3x+5x-4x-2+7=4x+5.
\]
}

\regola{Semplificare presto significa fare meno conti.}

% ------------------------------------------------------------
\voce{2. Eliminare presto i denominatori numerici}

\err{Portarsi dietro inutilmente le frazioni fino alla fine.}

Per esempio:
\[
\frac{x}{6}-\frac{x-1}{4}=\frac13.
\]

Il m.c.m. di $6,4,3$ è $12$.

\corr{Moltiplichiamo \textbf{entrambi i membri} per $12$:}
\[
2x-3(x-1)=4,
\]
\[
2x-3x+3=4,
\qquad
x=-1.
\]

\regola{Se possiamo liberarci facilmente delle frazioni,
conviene farlo subito.}

% ------------------------------------------------------------
\voce{3. Non usare due volte lo stesso fattore}

\err{
\[
2\cdot6(2x+1)
\longrightarrow
12(4x+2).
\]
}

Il fattore $2$ è stato utilizzato due volte.

\corr{
\[
2\cdot6(2x+1)=12(2x+1),
\]
oppure
\[
2\cdot6(2x+1)=6(4x+2).
\]
}

Entrambe danno:
\[
24x+12.
\]

% ------------------------------------------------------------
\voce{4. Segno meno davanti a una parentesi}

\err{
\[
-(2x-3+5x)=-2x-3+5x.
\]
}

\corr{
\[
-(2x-3+5x)=-2x+3-5x.
\]
}

Il meno equivale a moltiplicare \textbf{tutta la parentesi} per $-1$.

% ------------------------------------------------------------
\voce{5. Proprietà distributiva}

\err{
\[
3x(2x-5)=6x^2-5.
\]
}

\corr{
\[
3x(2x-5)
=6x^2-15x.
\]
}

Con due polinomi:
\[
(x+2)(x-3)
=x^2-3x+2x-6
=x^2-x-6.
\]

\regola{Ogni termine deve essere moltiplicato correttamente
per tutti i termini necessari.}

% ------------------------------------------------------------
\voce{6. Prodotti notevoli}

\err{
\[
(x-3)^2=x^2-9.
\]
}

\corr{
\[
(x-3)^2=x^2-6x+9.
\]
}

Invece:
\[
x^2-9=(x-3)(x+3)
\]
è una \textbf{differenza di quadrati}.

% ------------------------------------------------------------
\voce{7. Scomporre il trinomio speciale}

Per
\[
x^2+sx+p
\]
cerchiamo due numeri $a,b$ tali che

\[
a+b=s,
\qquad
ab=p.
\]

\err{Scambiare somma e prodotto, pensando che
\[
a+b=p,
\qquad
ab=s.
\]
}

Per esempio, davanti a
\[
x^2-5x+6
\]
cercare due numeri con prodotto $-5$ e somma $6$.

\corr{Bisogna invece cercare due numeri che abbiano}

\[
\text{somma}=-5,
\qquad
\text{prodotto}=6.
\]

Sono:
\[
-2,\,-3.
\]

Quindi:
\[
x^2-5x+6=(x-2)(x-3).
\]

\regola{Se hai un dubbio, \textbf{rimoltiplica} i fattori:}
\[
(x-2)(x-3)
=
x^2-3x-2x+6
=
x^2-5x+6.
\]

Se non riottieni il trinomio iniziale, la scomposizione è sbagliata.

% ------------------------------------------------------------
\voce{8. La legge di annullamento vale per il prodotto}

La legge di annullamento si usa quando abbiamo un
\textbf{prodotto uguale a zero}.

\[
(x-2)(x+5)=0
\]

significa:
\[
x-2=0
\qquad\text{oppure}\qquad
x+5=0.
\]

Quindi:
\[
S=\{-5;\,2\}.
\]

\regola{
\[
A\cdot B=0
\quad\Longrightarrow\quad
A=0 \ \text{oppure}\ B=0.
\]
}

% ------------------------------------------------------------
\voce{9. Non esiste la ``legge di annullamento della somma''}

\err{Da
\[
x^2-5x+6=0
\]
scrivere
\[
x^2=0,\qquad -5x=0,\qquad 6=0.
\]
}

Questo procedimento \textbf{non ha alcun significato}.

\corr{Prima bisogna scomporre:}
\[
x^2-5x+6
=
(x-2)(x-3).
\]

Poi:
\[
(x-2)(x-3)=0
\]
e quindi
\[
x=2
\qquad\text{oppure}\qquad
x=3.
\]

\[
S=\{2;\,3\}.
\]

\regola{Da $A+B=0$ non segue che $A=0$ e $B=0$.}

% ------------------------------------------------------------
\voce{10. ``Trasportare'' un termine}

\err{
\[
3x+5=11
\quad\Rightarrow\quad
3x=11+5.
\]
}

\corr{Sottraiamo $5$ da entrambi i membri:}
\[
3x+5-5=11-5,
\]
quindi
\[
3x=6.
\]

\regola{Il cambio di segno deriva dal fare la stessa
operazione sui due membri: non è una regola magica.}

\end{multicols}

% ============================================================
% PAGINA 2
% ============================================================

\newpage

\begin{center}
	{\large\bfseries Errori ricorrenti -- seconda parte}
\end{center}

\vspace{-2mm}

\begin{multicols}{2}
\small

% ------------------------------------------------------------
\voce{11. Grado dell'equazione e numero di soluzioni}

\err{Dire che
\[
(x-2)^2=0
\]
ha due soluzioni perché è di secondo grado.}

\corr{
\[
(x-2)(x-2)=0.
\]
}

Entrambi i fattori si annullano per lo stesso valore:
\[
x=2.
\]

Quindi:
\[
S=\{2\}.
\]

% ------------------------------------------------------------
\voce{12. Principi di equivalenza}

Non basta dire: ``faccio la stessa cosa ai due membri''.

\err{Da
\[
x^2=4x
\]
dividere immediatamente per $x$:
\[
x=4.
\]
}

Così si perde la possibile soluzione $x=0$.

\corr{
\[
x^2-4x=0,
\]
\[
x(x-4)=0,
\]
\[
x=0
\quad\text{oppure}\quad
x=4.
\]
}

Quindi:
\[
S=\{0;\,4\}.
\]

\regola{Prima di moltiplicare o dividere per un'espressione
contenente l'incognita bisogna chiedersi se può valere zero.}

% ------------------------------------------------------------
\voce{13. Dominio e insieme delle soluzioni}

Il \textbf{dominio} indica i valori per cui le espressioni hanno senso.

L'\textbf{insieme delle soluzioni} contiene i valori del dominio
che rendono vera l'equazione.

Per esempio:
\[
\frac1{x-2}=1.
\]

Dominio:
\[
D=\R\setminus\{2\}.
\]

Soluzione:
\[
x=3,
\qquad
S=\{3\}.
\]

% ------------------------------------------------------------
\voce{14. Scrivere correttamente l'insieme delle soluzioni}

\err{
\[
S=\{\varnothing\},
\qquad
S=\{\R\}.
\]
}

\corr{Se non ci sono soluzioni:}
\[
S=\varnothing.
\]

Se tutti i reali sono soluzioni:
\[
S=\R.
\]

Se le soluzioni sono $2$ e $5$:
\[
S=\{2;\,5\}.
\]

\regola{
\[
\varnothing\neq\{\varnothing\},
\qquad
\R\neq\{\R\}.
\]
}

Le parentesi graffe significano che ciò che è scritto dentro
è un \textbf{elemento} dell'insieme.

% ------------------------------------------------------------
\voce{15. Coefficienti frazionari ed equazioni frazionarie}

\err{Dire che
\[
\frac12x+\frac13=2
\]
è un'equazione frazionaria soltanto perché compaiono frazioni.}

\corr{È \textbf{intera rispetto a $x$}, perché l'incognita
non compare nei denominatori.}

Invece:
\[
\frac1{x-2}=3
\]
è frazionaria rispetto a $x$.

% ------------------------------------------------------------
\voce{16. Ricavare una grandezza da una formula}

Consideriamo:
\[
P=\frac{F}{S}.
\]

Vogliamo isolare $S$.

\err{Scambiare le lettere a intuito, per esempio:
\[
S=\frac{P}{F}.
\]
}

\corr{Moltiplichiamo entrambi i membri per $S$:}
\[
PS=F.
\]

Poi dividiamo entrambi i membri per $P$:
\[
S=\frac{F}{P}.
\]

\regola{Una formula è un'equazione:
si usano operazioni inverse su \textbf{entrambi i membri}.}

% ------------------------------------------------------------
\voce{17. ``Una frazione di una quantità''}

\err{Interpretare ``un quinto di $P$'' come
\[
P:\frac15.
\]
}

\corr{
\[
\frac15P=P:5.
\]
}

Per esempio:
\[
\frac15\cdot50=10.
\]

% ------------------------------------------------------------
\voce{18. Sconto e prezzo finale}

Se un oggetto costa $P$ e lo sconto è di un quinto:

\[
\text{sconto}=\frac15P.
\]

\err{Scrivere che il prezzo finale è
\[
\frac15P.
\]
}

\corr{Lo sconto deve essere sottratto:}
\[
P-\frac15P
=
\frac45P.
\]

Se il prezzo è $50$ euro:
\[
\frac15\cdot50=10
\]
è lo sconto,

\[
50-10=40
\]
è il prezzo pagato.

% ------------------------------------------------------------
\voce{19. Tradurre il linguaggio in matematica}

La matematica non consiste soltanto nel manipolare simboli.

Bisogna saper interpretare, per esempio:

\[
\text{``un quinto di }x\text{''}
\longrightarrow
\frac{x}{5},
\]

\[
\text{``il doppio di }x+3\text{''}
\longrightarrow
2(x+3),
\]

\[
\text{``la metà della differenza tra }x\text{ e }4\text{''}
\longrightarrow
\frac{x-4}{2}.
\]

% ------------------------------------------------------------
\voce{20. Prima di continuare, fermati e chiediti}

\begin{enumerate}[
	leftmargin=*,
	itemsep=0.5pt,
	topsep=1pt,
	parsep=0pt
]
	\item Posso \textbf{semplificare} qualcosa?
	\item Posso eliminare subito dei \textbf{denominatori numerici}?
	\item Ho sviluppato correttamente parentesi e prodotti?
	\item Posso o devo \textbf{scomporre}?
	\item Se ho scomposto, posso \textbf{rimoltiplicare per controllare}?
	\item Ho davvero un \textbf{prodotto uguale a zero} prima di usare la LAP?
	\item Sto applicando la stessa operazione ai \textbf{due membri}?
	\item Il passaggio conserva davvero le \textbf{soluzioni}?
	\item Quello che ho scritto ha \textbf{senso matematico}?
\end{enumerate}

\vspace{2mm}

\begin{center}
\fbox{
	\parbox{0.91\columnwidth}{
		\centering
		\textbf{L'obiettivo non è fare tanti passaggi.}\\
		È fare pochi passaggi, corretti, semplici e consapevoli.
	}
}
\end{center}

\end{multicols}

\end{document}
